% =============================================================================
%  Sequre — Bachelor's Thesis Documentation
%  Author : Vezeteu Andrei
%  Faculty: Alexandru Ioan Cuza University of Iasi, Faculty of Computer Science
%  Year   : 2025-2026
% =============================================================================

\documentclass[12pt, a4paper]{article}

% ---------- Encoding & fonts ----------
\usepackage[utf8]{inputenc}
\usepackage[T1]{fontenc}
\usepackage{lmodern}
\usepackage[english]{babel}
\usepackage{microtype}

% ---------- Layout ----------
\usepackage{geometry}
\geometry{
  top=2.5cm, bottom=2.5cm,
  left=3cm,  right=2.5cm
}
\usepackage{parskip}
\usepackage{setspace}
\onehalfspacing

% ---------- Section formatting ----------
\usepackage{titlesec}
\titleformat{\section}{\Large\bfseries}{\thesection}{1em}{}
\titleformat{\subsection}{\large\bfseries}{\thesubsection}{1em}{}
\titleformat{\subsubsection}{\normalsize\bfseries}{\thesubsubsection}{1em}{}

% ---------- Lists ----------
\usepackage{enumitem}
\setlist[itemize]{noitemsep, topsep=4pt}

% ---------- Headers & footers ----------
\usepackage{fancyhdr}
\pagestyle{fancy}
\fancyhf{}
\rhead{\small Sequre --- Bachelor's Thesis}
\lhead{\small Vezeteu Andrei}
\cfoot{\thepage}
\renewcommand{\headrulewidth}{0.4pt}

% ---------- Graphics, tables, math ----------
\usepackage{graphicx}
\usepackage{booktabs}
\usepackage{array}
\usepackage{xcolor}
\usepackage{amsmath, amssymb}
\usepackage{physics}
\usepackage{braket}
\usepackage{float}
\usepackage{caption}
\captionsetup[figure]{labelfont={small,it}, font=small, labelsep=endash}
\captionsetup[table]{labelfont={small,it}, font=small, labelsep=endash}

% ---------- Code listings ----------
\usepackage{listings}
\lstset{
  basicstyle=\ttfamily\footnotesize,
  breaklines=true,
  frame=single,
  captionpos=b,
  showstringspaces=false,
  numbers=left,
  numberstyle=\tiny\color{gray},
  keywordstyle=\color{blue!70!black}\bfseries,
  commentstyle=\color{green!50!black}\itshape,
  stringstyle=\color{orange!80!black}
}

% ---------- Hyperlinks ----------
\usepackage{hyperref}
\hypersetup{
  colorlinks=true,
  linkcolor=black,
  citecolor=blue!60!black,
  urlcolor=blue!70!black,
  pdftitle={Sequre --- Bachelor's Thesis},
  pdfauthor={Vezeteu Andrei}
}

% =============================================================================
\begin{document}

% =================== TITLE PAGE ==============================================
\begin{titlepage}
  \begin{center}
    {\large \textbf{``ALEXANDRU IOAN CUZA'' UNIVERSITY OF IA\c{S}I}} \\[4pt]
    {\large \textbf{FACULTY OF COMPUTER SCIENCE}} \\[40pt]

    \includegraphics[width=4cm]{figures/fii-logo.png} \\[30pt]

    {\Large BACHELOR'S THESIS} \\[10pt]
    {\LARGE \textbf{Sequre: A Real-Time Simulation and \\[4pt]
                    Visualization Platform for the BB84 \\[4pt]
                    Quantum Key Distribution Protocol}} \\[40pt]

    {\normalsize submitted by} \\[14pt]
    {\large \textbf{Vezeteu Andrei}} \\[60pt]

    \begin{minipage}{0.45\textwidth}
      \raggedright
      \textbf{Session:} July, 2026
    \end{minipage}\hfill
    \begin{minipage}{0.45\textwidth}
      \raggedleft
      \textit{Scientific coordinator} \\
      \textbf{Conf.\ Dr.\ Andreea Arusoaie}
    \end{minipage}

  \end{center}
\end{titlepage}

% =================== ABSTRACT ================================================
\newpage
\thispagestyle{plain}
\begin{center}
  {\Large \textbf{Abstract}}
\end{center}
\vspace{10pt}

\noindent
% TODO: 1-page abstract describing Sequre — to be filled in once chapters are stable.
\textit{[Abstract to be written after the main chapters are finalised, so the
summary accurately reflects the final content. Expected length: ~one page.
Topics to cover: the problem (need for accessible QKD experimentation),
the proposed solution (Sequre platform), key technical contributions (real
IBM Quantum hardware integration, real-time WebSocket visualization, full
end-to-end BB84 pipeline with EC and PA, multi-mode Eve simulation),
methodology, and main results.]}

\vspace{20pt}
\noindent
\textit{Keywords:} BB84, Quantum Key Distribution, QKD, Quantum Cryptography,
Qiskit, IBM Quantum, FastAPI, React, WebSockets, QBER, Privacy Amplification,
Error Correction, No-Cloning Theorem.

% =================== TABLE OF CONTENTS =======================================
\newpage
\tableofcontents
\thispagestyle{plain}

% =================== 1. INTRODUCTION =========================================
\newpage
\section{Introduction}
\label{sec:introduction}

\subsection{Motivation}
\label{sec:motivation}

The security of virtually all digital communication today rests on the
computational difficulty of certain mathematical problems --- factoring large
integers in the case of RSA, and computing discrete logarithms in the case of
elliptic-curve cryptography (ECC). These schemes are not proven secure in an
information-theoretic sense; they are secure only because no classical
polynomial-time algorithm is known to break them.

This foundation was shaken in 1994, when Peter Shor demonstrated that a
sufficiently large quantum computer could factor an $n$-bit integer in
polynomial time~\cite{shor1994}. Although practical fault-tolerant quantum
computers capable of breaking 2048-bit RSA remain several years away, the
cryptographic community has already entered a transition period. The
\textit{harvest now, decrypt later} strategy --- in which adversaries collect
encrypted traffic today with the intent to decrypt it once quantum hardware
matures --- represents a concrete, near-term threat to the confidentiality of
long-lived sensitive data~\cite{mosca2018}.

In response to this threat, the National Institute of Standards and Technology
(NIST) launched a post-quantum cryptography standardisation process in 2016,
selecting its first set of quantum-resistant algorithms in 2024. However,
\textit{post-quantum cryptography} (PQC) still relies on mathematical hardness
assumptions, merely choosing problems believed to be resistant to quantum
algorithms. A fundamentally different approach, one that achieves
\textit{information-theoretic} security whose guarantees are grounded in the
laws of physics rather than computational assumptions, is offered by
\textbf{Quantum Key Distribution} (QKD).

QKD allows two parties, commonly referred to as Alice and Bob, to establish a
shared secret key over an insecure channel such that any eavesdropping attempt
by a third party (Eve) is detectable with overwhelming probability. The
security follows from two foundational principles of quantum mechanics: the
\textbf{no-cloning theorem}, which forbids copying an arbitrary unknown quantum
state~\cite{wootters1982}, and the disturbance introduced by measurement in an
incompatible basis, as formalised by the
\textbf{Heisenberg uncertainty principle}. Any interception by Eve necessarily
disturbs the quantum states she measures, leaving a statistically detectable
signature in the error rate of the received key.

The first and most influential QKD protocol was proposed by Charles Bennett and
Gilles Brassard in 1984~\cite{bb84}. The BB84 protocol encodes classical bits
into the polarisation states of individual photons using two mutually unbiased
bases, and its unconditional security has been rigorously proved under realistic
noise models~\cite{mayers2001,loChau1999}. Despite its theoretical maturity,
however, BB84 remains inaccessible to most students and researchers for
practical experimentation: real QKD hardware is prohibitively expensive,
and existing software simulators are either too narrowly focused on performance
benchmarking, lack real-time visualisation, or do not integrate with actual
quantum processors.

This accessibility gap motivates the present work. There is a clear need for an
interactive, web-based platform that can demonstrate the full BB84 pipeline ---
from qubit generation through to a derived 128-bit secret key --- while offering
real hardware execution on IBM Quantum backends and providing a step-by-step
visual intuition for every stage of the protocol. Such a tool would serve both
an educational purpose, making quantum cryptography tangible for students
encountering it for the first time, and a research purpose, providing a
configurable testbed for studying protocol behaviour under varied noise regimes
and adversarial conditions.

\subsection{Proposed Solution}
\label{sec:proposed-solution}

\textbf{Sequre} is a real-time simulation and visualisation platform for the
BB84 Quantum Key Distribution protocol. It is implemented as a web application
accessible through any modern browser, with no local installation required
beyond a running Docker environment. The platform provides a complete,
end-to-end demonstration of the quantum cryptographic pipeline under three
distinct execution modes.

At its core, Sequre implements the full BB84 protocol in software using
\textbf{Qiskit}~\cite{qiskit2024} for quantum circuit construction and
\textbf{Qiskit-Aer} for noise-model simulation. Alice encodes classical bits
into single-qubit states drawn from the rectilinear basis ($\ket{0}$,
$\ket{1}$) or the diagonal basis ($\ket{+}$, $\ket{-}$), and Bob measures each
received qubit in a randomly chosen basis. Following the exchange, the sifting
step retains only those bits where Alice and Bob chose the same basis (on
average 50\% of transmitted qubits). The platform then computes the
\textbf{Quantum Bit Error Rate} (QBER), applies error correction via a
parity-block reconciliation scheme, and performs privacy amplification by
truncating a SHA-256 digest of the reconciled key to 128 bits.

The simulation operates in one of three modes selectable by the user:
\begin{itemize}
  \item \textbf{Ideal simulation:} noise-free quantum circuits on a statevector
        or QASM simulator, producing a zero-QBER baseline.
  \item \textbf{Noisy simulation:} configurable depolarising and measurement
        errors applied via Qiskit-Aer, producing realistic QBER distributions
        that allow the user to study the protocol's behaviour as a function of
        hardware imperfection.
  \item \textbf{Real IBM Quantum hardware:} circuits are transpiled and
        dispatched to a real superconducting quantum processor via
        \textbf{Qiskit IBM Runtime} (SamplerV2 API), currently supporting the
        \texttt{ibm\_fez}, \texttt{ibm\_marrakesh}, and
        \texttt{ibm\_kingston} backends. Results are polled asynchronously and
        streamed live to the frontend.
\end{itemize}

Beyond the core protocol, Sequre implements several layers of physical and
analytical realism that bring the simulator closer to a research-grade
testbed:

\begin{itemize}
  \item \textbf{Optical-fibre channel attenuation.} Each transmitted photon
        survives the channel with probability
        $\eta(L) = 10^{-\alpha L/10}$, where $\alpha = 0.2$\,dB/km is the
        standard single-mode-fibre loss coefficient and $L$ is the
        configurable transmission distance. A dedicated chart plots the
        asymptotic secure key rate $R \sim \eta(L)(1 - 2h(Q))$ across
        $0$--$200$\,km, illustrating the well-known rate--distance limits
        of point-to-point QKD without trusted nodes.

  \item \textbf{Decoy-state BB84 with weak coherent pulses.} The user can
        switch from an idealised single-photon source to a
        \emph{weak-coherent-pulse} (WCP) source emitting two intensity
        classes (signal $\mu_s$ and decoy $\mu_d$), with Poissonian
        photon-number statistics. Per-intensity gains and error rates are
        tracked online, and the GLLP two-decoy bounds are evaluated to
        estimate the single-photon yield $Y_1$ and error rate $e_1$ ---
        the quantities required for security against the
        photon-number-splitting (PNS) attack.

  \item \textbf{Full CASCADE error correction.} The original parity-block
        scheme was replaced with the multi-pass CASCADE algorithm with
        back-correction, exposing both the corrected key and the
        information leakage $\ell$ that is later consumed by privacy
        amplification.

  \item \textbf{Machine-learning eavesdropper detection.} A Random Forest
        classifier trained on six contextual features per run
        (QBER, sifting efficiency, depolarising noise, measurement-error
        rate, channel distance, protocol length) produces an estimated
        probability that an eavesdropper was present. The classifier is
        retrainable on demand from the SQLite simulation history and
        provides a security signal that operates orthogonally to the QBER
        threshold.

  \item \textbf{User accounts and per-user attribution of runs.} A full
        authentication layer (email/password with bcrypt, JWT-based
        sessions in HttpOnly cookies, CSRF double-submit protection,
        password reset, and Google OAuth 2.0 sign-in) allows simulation
        runs to be attributed to authenticated users while keeping
        anonymous access fully functional.
\end{itemize}

An optional \textbf{Eve eavesdropping simulation} implements three distinct
attack strategies of progressively increasing sophistication:
\textit{Weak Eve}, which intercepts 30\% of transmitted qubits;
\textit{Strong Eve}, which intercepts all qubits and reliably generates a
QBER of approximately 25\%, well above the 11\% security threshold; and
\textit{Smart Eve}, an adaptive intercept-resend adversary that uses a
proportional feedback loop on the running QBER to tune its interception
probability towards a user-configurable target (e.g.\ 9\%, just below the
abort threshold), thereby evading a naive QBER-only security check. The
Smart Eve mode motivates the machine-learning detector described above and
provides a concrete adversarial baseline against which the classifier is
evaluated.

The user interface, built with \textbf{React} (Vite) and \textbf{Tailwind CSS},
receives per-qubit events from the backend over a persistent WebSocket
connection and renders them live without page refresh. The \textbf{PhotonGrid}
component animates the first 80 qubit exchanges cell by cell, encoding each
cell's basis symbol, bit value, and Bob's measurement basis, and highlights
Eve-intercepted qubits with a visual marker. A \textbf{key distillation
pipeline} component visualises the progressive shrinkage of the key through
each post-processing stage, while the \textbf{QBER historical chart} plots
error rates across all previous simulation runs against the security threshold
line. An \textbf{educational panel} of five expandable cards introduces the
underlying theory in accessible language alongside the simulation results.

Completed simulation runs are automatically persisted to a \textbf{SQLite}
database, and REST endpoints allow history retrieval and selective deletion,
enabling longitudinal analysis of QBER behaviour across multiple experiments.

\subsection{Contributions and Impact}
\label{sec:impact}

The primary contributions of this work are the following:

\begin{enumerate}[label=\textbf{C\arabic*.}, leftmargin=*, itemsep=6pt]

  \item \textbf{A fully functional, end-to-end BB84 implementation on real
        quantum hardware.} Unlike purely classical simulators, Sequre executes
        single-qubit BB84 circuits on IBM superconducting processors via the
        Qiskit IBM Runtime SamplerV2 interface. This integration enables the
        direct observation of real decoherence effects --- gate errors,
        readout noise, crosstalk --- on the QKD error rate, a dimension that
        no software-only simulator can replicate.

  \item \textbf{Real-time, event-driven visualisation of the quantum protocol.}
        The WebSocket streaming architecture delivers per-qubit simulation
        events to the frontend as they are generated, producing a live,
        animated view of the key exchange. This level of temporal granularity
        is absent from existing QKD educational tools, which typically present
        only aggregate results after a simulation completes.

  \item \textbf{A spectrum of eavesdropping models, from passive to adaptive.}
        Beyond the textbook intercept-resend Eve, Sequre implements a
        \emph{Smart Eve} adversary that uses a proportional feedback loop on
        the running QBER to evade the standard abort threshold while still
        extracting partial information. Combined with the QBER-based
        security check, this stack of three Eve modes (Weak, Strong, Smart)
        demonstrates concretely both the security guarantees and the
        \emph{limitations} of a threshold-only check against an intelligent
        adversary.

  \item \textbf{A realistic optical-fibre channel model with rate--distance
        curves.} A configurable attenuation coefficient
        ($\alpha = 0.2$\,dB/km by default) is applied per qubit, and the
        asymptotic secure key rate is plotted as a function of channel
        length, bringing the simulator in line with standard QKD
        benchmarking methodology.

  \item \textbf{Decoy-state BB84 with weak-coherent-pulse modelling.} Two
        intensity classes (signal and decoy) are emitted with Poissonian
        photon-number statistics, and the GLLP two-decoy bounds are
        evaluated online to recover the single-photon yield and error rate
        --- the security-critical quantities under the PNS attack.

  \item \textbf{Full CASCADE error correction with information-leakage
        accounting.} The multi-pass binary-search reconciliation with
        back-correction replaces the lossy parity-block scheme and exposes
        the leakage $\ell$ that feeds into privacy amplification, closing
        the loop on a rigorous post-processing pipeline.

  \item \textbf{Supervised eavesdropper detection via machine learning.}
        A Random Forest classifier trained on six contextual features per
        run produces a probability of eavesdropper presence that is
        complementary to (and orthogonal in failure mode to) the QBER
        threshold, providing a defence against precisely the adaptive
        adversaries that the QBER check cannot catch on its own.

  \item \textbf{Multi-user platform with authenticated session management.}
        A full authentication layer (email/password with bcrypt, JWT
        access/refresh tokens in HttpOnly cookies, CSRF double-submit
        protection, password reset, Google OAuth 2.0) allows simulations to
        be attributed to authenticated users, paving the way for the
        upcoming gamified and multi-player modes while preserving anonymous
        guest access.

  \item \textbf{A configurable research testbed for QKD experimentation.}
        The platform exposes noise parameters (depolarising probability,
        measurement error rate, interception fraction, channel distance,
        source type, decoy intensities, error-correction algorithm) as
        user-controlled inputs, making it suitable not only for classroom
        demonstration but also as an exploratory tool for studying the
        protocol's behaviour under varied hardware and adversarial
        conditions.

  \item \textbf{Persistent simulation history with REST API access.}
        Automated SQLite persistence of every simulation run, queryable via
        REST endpoints, enables longitudinal analysis and reproducible
        comparison of results across hardware modes, noise levels, and Eve
        configurations, while also serving as the training corpus for the
        ML classifier.

  \item \textbf{Multi-protocol support: B92, SARG04, and E91 with CHSH Bell
        test.} Three additional QKD protocols are implemented alongside BB84.
        B92 uses only two non-orthogonal states and conclusive-outcome sifting
        ($\approx 25\%$ efficiency). SARG04 matches BB84's four states but
        announces \emph{pairs} of states at sifting time, providing resistance
        against the PNS attack without requiring a decoy source. E91 generates
        entangled Bell pairs ($\ket{\Phi^+}$), distributes one photon to each
        party, and derives the key from the correlated measurement outcomes;
        eavesdropping is certified by computing the CHSH correlator $S$
        and verifying $|S| > 2$, the classical bound. The platform visualises
        $S$ in real time and flags Bell-inequality violation or collapse.

  \item \textbf{Explicit Photon-Number-Splitting attack simulator.} A dedicated
        Eve mode models the PNS attack on WCP sources: single-photon pulses
        are blocked (no information obtainable without disturbing the state),
        while multi-photon pulses are split with one photon retained in quantum
        memory and the rest forwarded losslessly. The resulting collapse of the
        single-photon yield $Y_1 \to 0$ and the secure key rate $R \to 0$, as
        computed by the Lo-Ma-Chen decoy-state bounds, is displayed in the
        PNS Attack panel alongside the fraction of pulses Eve successfully
        intercepted.

  \item \textbf{Composable finite-key security bounds (Tomamichel 2012).}
        A dedicated analysis module computes the Hoeffding statistical
        correction $\delta_\mathrm{PE} = \sqrt{\ln(2/\varepsilon_\mathrm{PE})/(2n)}$
        on the QBER estimate, the Slepian--Wolf error-correction leakage, and
        the privacy-amplification overhead, yielding both the asymptotic
        (textbook GLLP) key length and the rigorous composable finite-key
        length for the actual block size $n$. The panel visualises the gap
        between the two and identifies the minimum $n$ required for a positive
        finite-key bound at a given security parameter $\varepsilon_\mathrm{sec}$.

  \item \textbf{LSTM deep-learning eavesdropper detector.} A recurrent neural
        network (PyTorch, hidden size 64, one LSTM layer) is trained on 1\,200
        synthetic per-qubit exchange sequences and learns to predict
        $P(\text{Eve} \mid \text{sequence})$ from the full temporal pattern of
        alice/bob basis choices, bob results, and error flags. This model
        achieves $92\%$ validation accuracy and correctly flags intercept-resend
        attacks at the operational 30\% interception rate --- a regime in which
        the QBER remains below the 11\% abort threshold and a threshold-only
        check would declare the channel secure. Model weights are cached on
        disk and loaded at startup; training completes in under 30 seconds on
        CPU.

\end{enumerate}

Beyond its immediate academic scope, Sequre is conceived as a
research-oriented testbed with a long-term ambition: to probe the practical
limits of quantum key distribution under realistic hardware imperfections,
advanced attack models, and extended protocol variants. The architecture is
deliberately modular --- the quantum logic, classical post-processing, and
presentation layers are cleanly separated --- and this modularity has already
enabled the integration of decoy-state BB84, CASCADE error correction, an
optical-fibre attenuation model, an adaptive (Smart) Eve, an ML-based
detection layer, three alternative QKD protocols (B92, SARG04, E91 with
CHSH Bell test), an explicit PNS-attack simulator, composable finite-key
security bounds, and a deep-learning (LSTM) eavesdropper detector into the
same code base without restructuring. A gamified \emph{Challenge Mode} is
identified as the next development target (see Section~\ref{sec:conclusions}).

% =================== 2. THEORETICAL BACKGROUND ===============================
\newpage
\section{Theoretical Background}
\label{sec:background}

This chapter provides the theoretical foundations necessary to understand the
Sequre platform and the BB84 protocol it implements. It begins with a
comparison of classical and quantum approaches to cryptographic key
distribution, introduces the relevant quantum mechanical formalism, and then
describes in detail the BB84 protocol, its security metric (QBER), and the
classical post-processing steps that transform raw qubit measurements into a
usable secret key.

\subsection{Classical vs.\ Quantum Cryptography}
\label{sec:classical-vs-quantum}

Classical cryptography can be divided into two broad categories:
\textit{symmetric-key} cryptography, in which sender and receiver share a
common secret key, and \textit{public-key} (asymmetric) cryptography, in which
a public key is used for encryption and a private key for decryption. The
security of public-key schemes such as RSA and elliptic-curve cryptography
(ECC) relies on the computational intractability of certain mathematical
problems --- integer factorisation and the discrete logarithm problem,
respectively. These problems are hard for classical computers, but as
established in Section~\ref{sec:motivation}, Shor's algorithm~\cite{shor1994}
solves both in polynomial time on a sufficiently powerful quantum computer.
Symmetric schemes are weakened but not broken: Grover's algorithm reduces an
$n$-bit key search from $O(2^n)$ to $O(2^{n/2})$~\cite{grover1996}, so
doubling the key length restores the original security margin.

\textbf{Quantum Key Distribution} takes an entirely different approach. Rather
than relying on mathematical assumptions, QKD uses the physical properties of
quantum states to distribute a key in such a way that any interception attempt
is detectable in principle. The security is \textit{information-theoretic}: it
holds regardless of the computational power of the adversary. The key insight
is that quantum mechanics imposes a fundamental constraint on any measurement
process --- measuring a quantum state in an incompatible basis irreversibly
disturbs that state, leaving a detectable trace.

Table~\ref{tab:crypto-comparison} summarises the key differences between
the two paradigms.

\begin{table}[H]
  \centering
  \caption{Classical vs.\ quantum key distribution --- a comparison.}
  \label{tab:crypto-comparison}
  \begin{tabular}{@{} l p{5.5cm} p{5.5cm} @{}}
    \toprule
    \textbf{Property} & \textbf{Classical (RSA/ECC)} & \textbf{QKD (BB84)} \\
    \midrule
    Security basis     & Computational hardness     & Laws of quantum physics \\
    Quantum threat     & Broken by Shor's algorithm & Unaffected \\
    Eavesdrop detection & Not inherent              & Guaranteed by physics \\
    Key distribution   & Public-key exchange        & Quantum channel + public classical channel \\
    Security model     & Computational              & Information-theoretic \\
    \bottomrule
  \end{tabular}
\end{table}

\subsection{Quantum Mechanics Primer}
\label{sec:qm-primer}

\subsubsection{Qubits and Quantum States}

The fundamental unit of quantum information is the \textbf{qubit} (quantum
bit). Unlike a classical bit, which is either 0 or 1, a qubit can exist in a
\textit{superposition} of both basis states simultaneously. Using Dirac
(bra-ket) notation, the general state of a qubit is written as:
\begin{equation}
  \ket{\psi} = \alpha\ket{0} + \beta\ket{1},
  \qquad \alpha,\beta \in \mathbb{C}, \quad |\alpha|^2 + |\beta|^2 = 1,
  \label{eq:qubit}
\end{equation}
where $\ket{0}$ and $\ket{1}$ are the computational basis states (eigenstates
of the Pauli-$Z$ operator), and $|\alpha|^2$, $|\beta|^2$ are the probabilities
of measuring the qubit in state $\ket{0}$ or $\ket{1}$, respectively.

\subsubsection{Measurement and State Collapse}

A key property of quantum mechanics is that measuring a qubit in a given basis
\textit{collapses} its state to one of the basis vectors. If the qubit is in
state $\ket{\psi} = \alpha\ket{0} + \beta\ket{1}$ and we measure in the
$Z$-basis $\{\ket{0},\ket{1}\}$, we obtain outcome $0$ with probability
$|\alpha|^2$ and outcome $1$ with probability $|\beta|^2$, and the state
collapses accordingly. After the measurement, all information about the
pre-measurement superposition is irrevocably lost.

Crucially, if a qubit is prepared in a state belonging to one orthonormal basis
and then measured in a \textit{different}, incompatible basis, the outcome is
uniformly random and the original state cannot be inferred. This is the
physical mechanism exploited by BB84 to detect eavesdropping.

\subsubsection{The Hadamard Gate and the Diagonal Basis}

In addition to the computational $Z$-basis $\{\ket{0},\ket{1}\}$, BB84 uses
the \textit{diagonal} ($X$-basis) $\{\ket{+},\ket{-}\}$, defined by:
\begin{equation}
  \ket{+} = \frac{\ket{0} + \ket{1}}{\sqrt{2}}, \qquad
  \ket{-} = \frac{\ket{0} - \ket{1}}{\sqrt{2}}.
  \label{eq:hadamard-states}
\end{equation}
The transformation between the two bases is performed by the
\textbf{Hadamard gate} $H$:
\begin{equation}
  H = \frac{1}{\sqrt{2}}\begin{pmatrix} 1 & 1 \\ 1 & -1 \end{pmatrix},
  \quad H\ket{0} = \ket{+}, \quad H\ket{1} = \ket{-}.
  \label{eq:hadamard}
\end{equation}
In the Sequre implementation, the Hadamard gate is applied in Qiskit via
\texttt{qc.h(0)} when Alice encodes a bit in the $X$-basis, and equally when
Bob measures in the $X$-basis (rotating back before the standard $Z$-measurement).

\subsection{The No-Cloning Theorem}
\label{sec:no-cloning}

\begin{quote}
  \textbf{Theorem} (Wootters and Zurek, 1982~\cite{wootters1982}): There is no
  unitary operation $U$ that can copy an arbitrary unknown quantum state.
  Formally, there exists no $U$ such that
  $U\ket{\psi}\ket{0} = \ket{\psi}\ket{\psi}$ for all $\ket{\psi}$.
\end{quote}

\textit{Proof sketch.} Suppose such a $U$ existed. Then for any two states
$\ket{\psi}$ and $\ket{\phi}$:
\begin{align}
  U(\ket{\psi}\ket{0}) &= \ket{\psi}\ket{\psi}, \label{eq:clone1}\\
  U(\ket{\phi}\ket{0}) &= \ket{\phi}\ket{\phi}. \label{eq:clone2}
\end{align}
Taking the inner product of both sides:
\begin{equation}
  \braket{\psi|\phi} = \braket{\psi|\phi}^2.
  \label{eq:clone-inner}
\end{equation}
This equation holds only if $\braket{\psi|\phi} = 0$ (orthogonal states) or
$\braket{\psi|\phi} = 1$ (identical states). For any other pair of states,
$0 < |\braket{\psi|\phi}| < 1$, so no such $U$ can exist. $\square$

\textbf{Implication for QKD.} The no-cloning theorem makes it impossible for
Eve to intercept a transmitted qubit, make a perfect copy, forward the original
to Bob, and measure the copy at her leisure. Any attempt by Eve to gain
information about the qubit requires her to interact with it, which --- as the
following section explains --- necessarily disturbs the quantum state.

\subsection{Heisenberg's Uncertainty Principle in QKD}
\label{sec:heisenberg}

Heisenberg's uncertainty principle states that certain pairs of physical
properties, known as \textit{complementary observables}, cannot both be known
simultaneously with arbitrary precision. In the context of BB84, the relevant
complementary observables are the Pauli-$Z$ and Pauli-$X$ operators,
corresponding to measurement in the computational basis and the diagonal basis,
respectively.

Consider a qubit sent by Alice in the $Z$-basis in state $\ket{0}$.
If Bob --- or Eve --- measures this qubit in the $X$-basis instead, the outcome
is uniformly random: $\ket{0} = \frac{1}{\sqrt{2}}(\ket{+}+\ket{-})$, so the
result is $+$ or $-$ with equal probability $\frac{1}{2}$. Moreover, after
this $X$-basis measurement, the qubit's state has collapsed to either $\ket{+}$
or $\ket{-}$ --- no longer $\ket{0}$. If this disturbed qubit is forwarded to
Bob and he measures in the $Z$-basis (as Alice intended), his result is again
uniformly random, with a 50\% probability of disagreeing with Alice's original
bit.

This is the precise mechanism by which eavesdropping is detected in BB84:
\begin{itemize}
  \item When Eve intercepts a qubit, she must choose a measurement basis without
        knowing which basis Alice used.
  \item She guesses the correct basis with probability $\frac{1}{2}$.
  \item When she guesses wrong (50\% of intercepts), her measurement disturbs
        the state.
  \item Of those disturbed qubits, Bob's measurement disagrees with Alice's
        original bit with probability $\frac{1}{2}$.
\end{itemize}
Each intercepted qubit therefore introduces an error with probability
$\frac{1}{2} \times \frac{1}{2} = \frac{1}{4}$. A complete intercept-resend
attack (Strong Eve) thus raises the QBER to approximately 25\%.

\subsection{The BB84 Protocol}
\label{sec:bb84}

The BB84 protocol~\cite{bb84} proceeds as follows.

\paragraph{Step 1 --- Qubit preparation (Alice).}
Alice generates a random bit string $\mathbf{a} = (a_1, \ldots, a_n)$ and an
independent random basis string $\mathbf{x} = (x_1, \ldots, x_n)$, where each
$x_i \in \{Z, X\}$. She encodes each bit $a_i$ into a qubit according to
Table~\ref{tab:bb84-states} and transmits the qubits one by one over the
quantum channel.

\begin{table}[H]
  \centering
  \caption{BB84 encoding: the four qubit states used in the protocol.}
  \label{tab:bb84-states}
  \begin{tabular}{@{} cccl @{}}
    \toprule
    \textbf{Bit} & \textbf{Basis} & \textbf{State} & \textbf{Description} \\
    \midrule
    0 & $Z$ (rectilinear) & $\ket{0}$ & spin-up / horizontal polarisation \\
    1 & $Z$ (rectilinear) & $\ket{1}$ & spin-down / vertical polarisation \\
    0 & $X$ (diagonal)    & $\ket{+} = \tfrac{1}{\sqrt{2}}(\ket{0}+\ket{1})$ & $+45°$ polarisation \\
    1 & $X$ (diagonal)    & $\ket{-} = \tfrac{1}{\sqrt{2}}(\ket{0}-\ket{1})$ & $-45°$ polarisation \\
    \bottomrule
  \end{tabular}
\end{table}

\paragraph{Step 2 --- Qubit measurement (Bob).}
Bob independently chooses a random basis string $\mathbf{y} = (y_1, \ldots, y_n)$
and measures each received qubit in the corresponding basis, recording his
outcomes $\mathbf{b} = (b_1, \ldots, b_n)$.

\paragraph{Step 3 --- Basis sifting.}
Alice and Bob communicate their basis choices $\mathbf{x}$ and $\mathbf{y}$
over a public classical channel (no bit values are revealed). They retain
only the positions where $x_i = y_i$ --- on average, 50\% of all transmitted
qubits. The retained bits form the \textit{sifted key}. In Sequre, this step
is implemented in \texttt{classical\_processing/sifting.py}:
\begin{lstlisting}[language=Python, caption={Basis sifting in Sequre.}]
def sift_keys(alice_bits, alice_bases, bob_bits, bob_bases):
    alice_key, bob_key = [], []
    for a_bit, a_basis, b_bit, b_basis in zip(
            alice_bits, alice_bases, bob_bits, bob_bases):
        if a_basis == b_basis:
            alice_key.append(a_bit)
            bob_key.append(b_bit)
    return alice_key, bob_key
\end{lstlisting}

\paragraph{Step 4 --- Error rate estimation (QBER).}
Alice and Bob sacrifice a random subset of their sifted key to estimate the
Quantum Bit Error Rate (QBER). If QBER exceeds a security threshold, they
abort the protocol (see Section~\ref{sec:qber}).

\paragraph{Step 5 --- Error correction.}
Alice and Bob apply an error-correction protocol over the public channel to
reconcile the remaining bits of their sifted keys
(see Section~\ref{sec:ec-pa}).

\paragraph{Step 6 --- Privacy amplification.}
To account for any partial information Eve may have obtained, Alice and Bob
apply a universal hash function to compress their reconciled key into a shorter
but statistically independent final key
(see Section~\ref{sec:ec-pa}).

\subsection{Quantum Bit Error Rate (QBER)}
\label{sec:qber}

The \textbf{Quantum Bit Error Rate} (QBER) is defined as the fraction of
sifted key bits on which Alice and Bob disagree:
\begin{equation}
  Q = \frac{1}{N} \sum_{i=1}^{N} \mathbf{1}[a_i \neq b_i],
  \label{eq:qber}
\end{equation}
where $N$ is the length of the sifted key and $a_i$, $b_i$ are the
corresponding bits held by Alice and Bob after sifting.

In a noise-free channel without eavesdropping, $Q = 0$. Channel noise
(depolarising errors, measurement errors) raises $Q$ by a small, bounded
amount. An active intercept-resend attack raises $Q$ significantly above any
noise floor. Specifically:

\begin{itemize}
  \item \textbf{No Eve, no noise:} $Q \approx 0$.
  \item \textbf{Weak Eve} (intercepts fraction $p$ of qubits): each intercepted
        qubit introduces an error with probability $\frac{1}{4}$ (wrong basis
        with probability $\frac{1}{2}$, error given wrong basis with probability
        $\frac{1}{2}$), so $Q \approx \frac{p}{4}$. For Sequre's Weak Eve
        ($p = 0.30$), this gives $Q \approx 7.5\%$.
  \item \textbf{Strong Eve} ($p = 1.0$, full intercept-resend): $Q \approx 25\%$.
\end{itemize}

Sequre uses a security threshold of $Q_{\text{th}} = 11\%$, consistent with
the standard BB84 security analysis~\cite{shorPreskill2000}: if $Q \geq 11\%$,
the protocol is aborted and no key is output. This threshold ensures that any
remaining information Eve holds about the key can be made exponentially small
by privacy amplification.

\subsection{Error Correction and Privacy Amplification}
\label{sec:ec-pa}

\subsubsection{Error Correction}

Even in the absence of eavesdropping, channel noise causes a non-zero QBER.
Before a key can be used, Alice and Bob must reconcile their sifted keys so
that they hold identical bit strings. This is performed via an
\textit{information reconciliation} protocol over the public classical channel.

Sequre implements a simplified parity-block scheme: the sifted key is divided
into consecutive blocks of $B = 4$ bits. For each block, Alice announces the
block's parity (XOR of its bits) over the public channel. Bob computes the
parity of his corresponding block. If the parities agree, both parties retain
the block. If they disagree (indicating at least one error within the block),
the entire block is discarded by both parties. The implementation in
\texttt{classical\_processing/error\_correction.py} is:

\begin{lstlisting}[language=Python, caption={Parity-block error correction in Sequre.}]
def correct_errors(alice_key, bob_key, block_size=4):
    corrected_alice, corrected_bob = [], []
    for i in range(0, len(alice_key) - block_size + 1, block_size):
        a_block = alice_key[i:i + block_size]
        b_block = bob_key[i:i + block_size]
        if parity(a_block) == parity(b_block):
            corrected_alice.extend(a_block)
            corrected_bob.extend(b_block)
    return corrected_alice, corrected_bob
\end{lstlisting}

This approach sacrifices key material (entire blocks are discarded rather than
individual erroneous bits), but guarantees that all retained blocks are
error-free with high probability for low to moderate QBER. The full
\textbf{CASCADE} algorithm~\cite{brassard1994} improves upon this by
iteratively narrowing down the position of individual erroneous bits through
binary search, achieving near-optimal reconciliation efficiency. CASCADE is
identified as a future development direction for Sequre
(see Section~\ref{sec:conclusions}).

\subsubsection{Privacy Amplification}

After error correction, Alice and Bob share an identical key, but Eve may hold
partial information about it: she observed parity announcements over the public
channel during error correction, and she may have obtained information from any
qubits she intercepted before the QBER check. Privacy amplification
\textit{compresses} the reconciled key using a hash function, reducing Eve's
information about the output key to a negligible quantity regardless of how
much she learned about the input~\cite{bennett1995pa}.

Sequre implements privacy amplification by computing the \textbf{SHA-256}
hash of the reconciled key and truncating the digest to 128 bits
(32 hexadecimal characters):

\begin{lstlisting}[language=Python, caption={SHA-256 privacy amplification in Sequre.}]
def amplify_privacy(key, output_bits=128):
    byte_length = (len(key) + 7) // 8
    key_int = int("".join(str(b) for b in key), 2) if key else 0
    key_bytes = key_int.to_bytes(byte_length, byteorder="big")
    digest = hashlib.sha256(key_bytes).hexdigest()
    return digest[: output_bits // 4]  # 32 hex chars = 128 bits
\end{lstlisting}

SHA-256 is a collision-resistant, one-way function and serves as a practical
approximation to the universal hash families used in the information-theoretic
proofs of privacy amplification. While not strictly optimal from a
universal-hashing standpoint --- optimal constructions use seeded hash families
whose seed is shared publicly --- SHA-256 truncation is a common and accepted
choice for prototype implementations, and provides 128 bits of output security
in the absence of attacks on the hash function itself.

\subsection{Optical-Fibre Channel Attenuation}
\label{sec:channel-loss}

The discussion so far has treated the quantum channel as a passive conduit
that may add gate-level noise but transmits every qubit. In any physical
implementation of QKD, however, the dominant cause of bit loss is not noise
but \emph{photon attenuation} along the channel. For optical fibre links,
the survival probability of a photon transmitted over a distance $L$
(in km) follows the Beer--Lambert law:
\begin{equation}
  \eta(L) = 10^{-\alpha L / 10},
  \label{eq:fibre-attenuation}
\end{equation}
where $\alpha$ is the fibre's loss coefficient. The de-facto industry value
for telecom-grade single-mode fibre at the $1550$\,nm wavelength is
$\alpha \approx 0.2$\,dB/km~\cite{ITU-T-G652}. This implies, for instance,
that only $\eta(50) \approx 10\%$ of photons survive a $50$-km link and
$\eta(100) \approx 1\%$ survive $100$\,km. Free-space links incur instead
\emph{beam-divergence} loss, modelled as $\eta \propto 1/L^2$ but with
substantially smaller coefficients per unit distance over short ranges.

The practical consequence for QKD is that the \emph{secure key rate} ---
the number of secret-key bits produced per channel use --- decays
exponentially with distance. In the asymptotic (large-$N$) regime, the
secret-key rate for BB84 with one-way classical communication is given by
the GLLP--Shor--Preskill formula~\cite{shorPreskill2000,gllp2004}:
\begin{equation}
  R(L) \;\geq\; \eta(L) \,\bigl[1 \;-\; H_2(e_b) \;-\; f(e_b)\,H_2(e_b)\bigr],
  \label{eq:gllp-rate}
\end{equation}
where $e_b$ is the bit-error rate, $H_2$ is the binary entropy
$H_2(x) = -x\log_2 x - (1-x)\log_2(1-x)$, and $f(e_b) \geq 1$ is the
efficiency factor of the error-correction algorithm relative to the
Shannon limit. The term $1 - H_2(e_b)$ accounts for the information Eve
could have gained from the quantum channel; the term $f(e_b) H_2(e_b)$
discounts the bits revealed during classical reconciliation.

Equation~\eqref{eq:gllp-rate} sets the well-known
\emph{rate--distance trade-off} of BB84: there exists a maximum distance
$L^*$ beyond which the right-hand side becomes negative and no positive
key rate can be obtained. For $\alpha = 0.2$\,dB/km, $e_b = 1\%$, and
CASCADE-typical $f \approx 1.16$, $L^* \approx 150$--$200$\,km. Surpassing
this limit requires either trusted intermediate nodes, twin-field
protocols~\cite{lucamarini2018}, or quantum repeaters --- a key
observation that motivates the Sequre rate-vs-distance visualisation in
Section~\ref{sec:implementation}.

\subsection{Decoy-State Method and the Photon-Number-Splitting Attack}
\label{sec:decoy-state}

The BB84 protocol as originally formulated assumes Alice transmits a
\emph{single-photon source} (SPS), so that any attempt by Eve to extract
information necessarily perturbs the qubit. In practice, true single-photon
sources are demanding to engineer; instead, the overwhelming majority of
deployed QKD systems use attenuated lasers --- so-called
\emph{weak coherent pulses} (WCP). A WCP with mean photon number $\mu$
emits, per pulse, $n$ photons with probability:
\begin{equation}
  P(n \mid \mu) \;=\; e^{-\mu}\,\frac{\mu^n}{n!}.
  \label{eq:poisson}
\end{equation}
For $\mu < 1$, the dominant contributions are $n = 0$ (vacuum) and $n = 1$
(single photon), but a non-negligible fraction of pulses contains $n \geq 2$
photons. These multi-photon pulses expose the protocol to a
\textbf{photon-number-splitting} (PNS) attack: Eve installs a quantum
non-demolition photon-number measurement at the output of Alice, blocks all
$n = 1$ pulses, and for every $n \geq 2$ pulse keeps one photon in a
quantum memory while forwarding the rest to Bob unaltered. After Alice and
Bob publicly announce their bases, Eve measures her stored photons in the
correct basis, obtaining the bit value with no QBER penalty whatsoever.

The PNS attack is fatal to naive WCP-BB84 over high-loss channels, because
in the limit $\eta(L) \to 0$ the legitimate single-photon signal can be
completely suppressed in favour of intercepted multi-photon contributions
without raising the QBER above the security threshold. The
\textbf{decoy-state method}~\cite{lo2005decoy,wang2005decoy} elegantly
defeats this attack by having Alice randomly switch among several intensity
classes:
\begin{itemize}
  \item a \emph{signal} intensity $\mu_s$ used for key generation;
  \item one or more \emph{decoy} intensities $\mu_d, \nu, \ldots$ used only
        for monitoring.
\end{itemize}
Because Eve does not know in advance the intensity of any individual
pulse, she cannot distinguish signal multi-photon pulses from decoy
multi-photon pulses. The per-intensity gains
$Q_\mu = \sum_n P(n \mid \mu) Y_n$ and error rates $E_\mu$ provide a
system of linear inequalities that bound the \emph{single-photon yield}
$Y_1$ and \emph{single-photon error rate} $e_1$ in a verifiable way. In
the two-decoy GLLP analysis, the lower bound on $Y_1$ is:
\begin{equation}
  Y_1^{L} \;\geq\;
    \frac{\mu}{\mu\nu - \nu^2}\!\left(
      Q_\nu e^{\nu} - Q_\mu e^{\mu}\,\frac{\nu^2}{\mu^2}
      - \frac{\mu^2 - \nu^2}{\mu^2}\,Y_0
    \right),
  \label{eq:y1-bound}
\end{equation}
where $\mu$ and $\nu$ are the signal and decoy means and $Y_0$ is the
background (dark-count) rate. The secure key rate is then:
\begin{equation}
  R \;\geq\; q\,\bigl[
    -Q_\mu f(E_\mu) H_2(E_\mu)
    + Q_1^{L}\bigl(1 - H_2(e_1^{U})\bigr)
  \bigr],
  \label{eq:decoy-rate}
\end{equation}
with $Q_1^L = \mu e^{-\mu} Y_1^L$. The Sequre decoy-state panel renders
each of these quantities live, so the student can watch the
single-photon contribution to the key being recovered from the otherwise
aggregate WCP statistics.

\subsection{The CASCADE Information-Reconciliation Algorithm}
\label{sec:cascade}

The block-parity reconciliation described in Section~\ref{sec:ec-pa} is a
useful pedagogical baseline but is highly suboptimal: it discards entire
blocks on any disagreement, leaking only $\log_2(B)$ bits less than the
optimum per corrected error but losing $B$ bits of key in the process.
The CASCADE algorithm~\cite{brassard1994}, introduced by Brassard and
Salvail in 1993, remains the de-facto reference for QKD error correction
to this day because it isolates and corrects individual erroneous bits
rather than discarding their surrounding blocks.

CASCADE proceeds in $k$ \emph{passes}, each pass parameterised by a
block size $B_i$. The first block size is chosen as a function of the
estimated QBER $Q$:
\begin{equation}
  B_1 \;\approx\; \left\lceil \frac{0.73}{Q} \right\rceil,
  \qquad B_{i+1} = 2\,B_i.
  \label{eq:cascade-blocksize}
\end{equation}
Within a pass, the key is permuted by a fresh random permutation $\pi_i$
and partitioned into consecutive blocks of size $B_i$. Alice announces
the parity of each block over the public channel; for every block whose
parity differs from Bob's, the parties perform a binary search ---
recursively bisecting the block and comparing parities of the halves ---
until a single erroneous bit is located and flipped.

The crucial element of CASCADE is its \textbf{back-correction} mechanism:
flipping a bit in pass $i$ changes the parities of every block in passes
$1, \ldots, i-1$ that contained that bit. A flipped bit in a previously
``correct'' block must have masked a second error. CASCADE therefore
revisits all earlier blocks whose parity has now changed and continues
binary-searching for the remaining errors. This cross-pass propagation
allows the algorithm to converge with high probability after only four to
six passes, leaking
\begin{equation}
  \ell \;\approx\; N \cdot f(Q) \cdot H_2(Q)
  \label{eq:cascade-leakage}
\end{equation}
bits over the public channel, with $f(Q) \approx 1.05$--$1.22$ for
typical QKD QBER values --- close to the Shannon lower bound of
$f = 1$~\cite{martinez2015cascade}. The leakage $\ell$ is precisely the
quantity subtracted during privacy amplification to obtain the final
secret-key length.

In Sequre, CASCADE is implemented in
\texttt{backend/classical\_processing/cascade.py} and replaces the
block-parity scheme. The legacy scheme remains available behind a UI
selector for didactic comparison: a side-by-side view shows how many
bits are wasted by block discarding versus how many leak through CASCADE
parities at the same input QBER.

\subsection{Machine Learning for QKD Security Monitoring}
\label{sec:ml-detection}
\label{sec:lstm-impl}

A QBER-threshold check is a powerful but coarse instrument. It catches
any adversary that exceeds the abort threshold but is, by construction,
blind to attacks calibrated to stay below it --- attacks that Sequre's
Smart Eve mode produces by design. A complementary, statistics-driven
detector is therefore desirable.

The detection of subtle adversarial behaviour from an experiment's
contextual signature (QBER, sifting efficiency, noise parameters,
distance, sample size) is a textbook \emph{binary classification}
problem: given a feature vector $\mathbf{x} \in \mathbb{R}^d$, predict
the label $y \in \{0, 1\}$ (\emph{Eve absent}, \emph{Eve present}). Among
classical classifiers, the \textbf{Random Forest}~\cite{breiman2001} is
particularly well-suited to this setting: it is a non-linear ensemble of
decision trees that handles feature interactions natively, is robust to
the heterogeneous, non-Gaussian distributions characteristic of QKD
features, exposes interpretable feature importances, and trains in
seconds on the dataset sizes ($10^2$--$10^3$ runs) typical of an
educational platform.

Concretely, the model computes
\begin{equation}
  \hat{P}(y = 1 \mid \mathbf{x}) \;=\;
  \frac{1}{T} \sum_{t=1}^{T} \hat{P}_t(y = 1 \mid \mathbf{x}),
  \label{eq:rf-prob}
\end{equation}
where each $\hat{P}_t$ is the leaf-frequency estimate of the $t$-th tree
trained on a bootstrap sample of the historical runs. The thesis uses
$T = 100$ trees with the scikit-learn default depth limits. Labels are
derived automatically from the configured \texttt{eve\_mode} of each
historical run: anything other than \texttt{eve\_mode == 'none'} is
treated as a positive instance.

While the Random Forest operates on six aggregate features per run, the
literature has demonstrated that deeper models can exploit the full
temporal structure of the qubit exchange. LSTM networks, which process
the per-qubit sequence as a time series, have been shown to reach
$>90\%$ accuracy against adaptive attackers in
CV-QKD~\cite{xu2024dlqkd}. Sequre implements both approaches:
the Random Forest for fast, interpretable per-run scoring and the LSTM
detector described below for sequence-level detection of subtle patterns
that summary statistics flatten away.

\subsubsection{LSTM Per-Qubit Sequence Detector}
\label{sec:lstm-impl}

The Random Forest classifier described above is trained once on the
accumulated simulation history and operates on six hand-engineered
aggregate features. While effective against obvious eavesdropping, it
cannot exploit the \emph{temporal ordering} of qubit-level events --- a
dimension that is informative under intercept-resend attacks, where
errors cluster on basis-matched positions where Eve chose the wrong
basis, rather than appearing uniformly at random.

To address this, Sequre includes a \textbf{Long Short-Term Memory (LSTM)}
classifier~\cite{hochreiter1997} that reads the entire per-qubit exchange
as an ordered sequence. Each qubit $i$ produces a six-dimensional feature
vector:
\begin{equation}
  \mathbf{x}_i = \bigl(
    a_i^\text{basis},\;
    b_i^\text{basis},\;
    a_i^\text{bit},\;
    b_i^\text{result},\;
    \mathbf{1}[a_i^\text{basis} = b_i^\text{basis}],\;
    \mathbf{1}[b_i^\text{result} \neq a_i^\text{bit}]
    \cdot \mathbf{1}[a_i^\text{basis} = b_i^\text{basis}]
  \bigr),
  \label{eq:lstm-features}
\end{equation}
where the last two components encode the basis-match indicator and the
error indicator restricted to basis-matched positions, respectively.
These are the positions Eve's intercept-resend attack contaminates.

The network architecture is compact by design --- CPU inference is
required during the live simulation stream:
\begin{itemize}
  \item \textbf{Input:} variable-length sequence $[\mathbf{x}_1, \ldots,
        \mathbf{x}_n] \in \mathbb{R}^{n \times 6}$
  \item \textbf{LSTM layer:} hidden size 64, one layer, batch-first
  \item \textbf{Pooling:} element-wise sum of the mean-pooled hidden
        states and the final hidden state $h_n$, yielding a 64-dimensional
        representation
  \item \textbf{Fully connected:} $64 \to 2$ (logits for
        \emph{secure} / \emph{compromised})
\end{itemize}

\paragraph{Synthetic training data.} Because real QKD runs accumulate
slowly during a session, the model is trained offline on a
\emph{synthesised} corpus of 1\,200 sequences. Each sequence is generated
by rolling the per-qubit BB84 exchange stochastically: for Eve-present
samples, the interception rate is drawn from $[0.25, 1.0]$ with a bias
towards the boundary regime $[0.25, 0.45]$ --- the hardest cases for
a threshold-only check. For Eve-absent samples, the channel noise is
drawn uniformly from $[0.5\%, 5\%]$. An additional 200 sequences are
held out for validation.

\paragraph{Training and results.} The model is trained for 60 epochs with
the Adam optimiser ($\eta = 2 \times 10^{-3}$) and cross-entropy loss.
On the held-out validation set the model achieves:
\begin{itemize}
  \item \textbf{Validation accuracy:} 92\%
  \item \textbf{Validation loss:} 0.181
  \item \textbf{Training accuracy:} 89.5\%
\end{itemize}
Critically, at the operational Weak Eve rate of 30\% interception (which
produces QBER $\approx 7.5\%$, below the 11\% abort threshold), the LSTM
reports $P(\text{Eve}) > 0.97$ consistently across repeated trials,
whereas the QBER check alone would declare the channel secure.

Weights are persisted to \texttt{backend/ml/lstm\_weights.pt} and
loaded at server startup; if no cached weights exist, training runs
automatically in the lifespan hook and completes in under 30 seconds on
CPU.

\subsection{Alternative QKD Protocols}
\label{sec:alt-protocols}

BB84 is the reference QKD protocol, but several variants have been
proposed to address different hardware constraints and threat models.
Sequre implements three alternatives.

\subsubsection{B92 (Bennett, 1992)}
\label{sec:b92}

B92~\cite{b92} simplifies BB84 by using only two non-orthogonal states:
Alice sends $\ket{0}$ to encode bit 0 and $\ket{+}$ to encode bit 1.
Bob measures randomly in the $Z$- or $X$-basis. A measurement is
\emph{conclusive} only if Bob's outcome is inconsistent with
the state Alice \emph{could} have sent: outcome $\ket{1}$ (from a $Z$
measurement) rules out $\ket{0}$, so Alice must have sent $\ket{+}$,
meaning bit 1; outcome $\ket{-}$ (from an $X$ measurement) rules out
$\ket{+}$, meaning bit 0. All other outcomes are discarded. The sifting
efficiency is approximately 25\%, lower than BB84's 50\%, and the
protocol is more sensitive to intercept-resend attacks because Eve's
basis ambiguity is larger.

\subsubsection{SARG04 (Scarani et al., 2004)}
\label{sec:sarg04}

SARG04~\cite{sarg04} uses the same four qubit states as BB84 but modifies
the public announcement step: instead of revealing her basis, Alice
announces an \emph{unordered pair} $\{s_0, s_1\}$ of non-orthogonal
states, one of which is the state she actually sent. Bob retains the bit
only if his measurement outcome is consistent with exactly one of the two
announced states. Because Eve cannot identify which of the two states was
sent without measuring, the protocol resists the PNS attack more
effectively than BB84 at the same intensity --- at the cost of a
$\approx 25\%$ sifting efficiency, matching B92.

\subsubsection{E91 (Ekert, 1991) and the CHSH Bell Test}
\label{sec:e91}

The E91 protocol~\cite{ekert1991} takes an entirely different approach:
security is grounded in quantum entanglement rather than the no-cloning
theorem alone. A source emits pairs of photons in the maximally entangled
Bell state:
\begin{equation}
  \ket{\Phi^+} = \frac{1}{\sqrt{2}}\bigl(\ket{00} + \ket{11}\bigr).
  \label{eq:bell-state}
\end{equation}
Alice and Bob each receive one photon and independently choose a
measurement angle from a canonical set:
\begin{align}
  \text{Alice:} &\quad \theta_a \in \{0°,\,22.5°,\,45°\}, \\
  \text{Bob:}   &\quad \theta_b \in \{22.5°,\,45°,\,67.5°\}.
  \label{eq:e91-angles}
\end{align}
Rounds where both parties chose the same angle (both at $22.5°$ or both
at $45°$) contribute to the \emph{key}, exploiting the perfect
anti-correlation of $\ket{\Phi^+}$ measurements in the same basis. The
remaining rounds are used to evaluate the CHSH correlator:
\begin{equation}
  S = E(0°,\,22.5°) - E(0°,\,67.5°) + E(45°,\,22.5°) + E(45°,\,67.5°),
  \label{eq:chsh}
\end{equation}
where $E(\theta_a, \theta_b) = P(\text{same}) - P(\text{different})$ is
the correlation coefficient. For a maximally entangled state,
$|S| = 2\sqrt{2} \approx 2.828$ (Tsirelson bound); classical
local-hidden-variable theories are bounded by $|S| \leq 2$. If
eavesdropping destroys the entanglement, the measured $|S|$ drops toward
the classical bound and the parties abort the key exchange. In Sequre's
implementation, $S$ is computed after every E91 run and visualised
alongside the classical and quantum bounds; an intercept-resend attack
(Strong Eve) collapses $S$ to $\approx 0$.

\subsection{Finite-Key Security Bounds}
\label{sec:finite-key-theory}

The textbook BB84 security proof assumes the number of exchanged qubits
$N \to \infty$. In practice, any real-world run uses a finite block, and
the QBER estimate $\hat{Q}$ computed from $n$ sifted bits is a random
variable with statistical fluctuations of order $1/\sqrt{n}$. A rigorous
\emph{composable} security statement must account for this uncertainty.

Tomamichel et al.~\cite{tomamichel2012} derive a tight finite-key bound
using a Hoeffding concentration inequality on the phase-error estimate.
The corrected QBER worst-case is:
\begin{equation}
  Q_\mathrm{upper} = \hat{Q} + \delta_\mathrm{PE}, \qquad
  \delta_\mathrm{PE} = \sqrt{\frac{\ln(2/\varepsilon_\mathrm{PE})}{2n}},
  \label{eq:hoeffding}
\end{equation}
where $\varepsilon_\mathrm{PE}$ is the parameter-estimation security
parameter. The composable finite-key length is then:
\begin{equation}
  \ell \leq n\bigl[1 - H_2(Q_\mathrm{upper})\bigr]
           - n\,f_\mathrm{EC}\,H_2(\hat{Q})
           - \log_2\!\tfrac{2}{\varepsilon_\mathrm{corr}}
           - 2\log_2\!\tfrac{1}{2\varepsilon_\mathrm{PA}},
  \label{eq:finite-key}
\end{equation}
where the second term is the error-correction leakage (Slepian--Wolf
bound with efficiency factor $f_\mathrm{EC} \approx 1.16$ for CASCADE)
and the final two terms are the privacy-amplification and correctness
overheads. The total security parameter is
$\varepsilon_\mathrm{sec} = \varepsilon_\mathrm{PE}
+ \varepsilon_\mathrm{corr} + \varepsilon_\mathrm{PA}$.

Equation~\eqref{eq:finite-key} has a characteristic behaviour: $\delta_\mathrm{PE}
\propto 1/\sqrt{n}$ shrinks as $n$ grows, so the finite-key length
converges from below to the asymptotic GLLP value. For typical security
parameters ($\varepsilon_\mathrm{sec} = 3 \times 10^{-10}$), the PA and
correctness overhead is a fixed cost of $\approx 100$ bits, which
dominates for small $n$ and becomes negligible only above
$n \approx 500$--$1\,000$ sifted bits. Sequre's Finite-Key panel
visualises this gap for the current run and identifies whether the block
size is sufficient for a positive bound.

% =================== 3. EXISTING QKD SIMULATORS ==============================
\newpage
\section{Existing QKD Simulators}
\label{sec:related-work}

% TODO: Survey of comparable platforms + how Sequre differs.

\subsection{QKDSim}
\subsection{NetSquid}
\subsection{QuNetSim / SimulaQron}
\subsection{Positioning of Sequre}

% =================== 4. TECHNOLOGIES USED ====================================
\newpage
\section{Technologies Used}
\label{sec:technologies}

\subsection{Python and Qiskit}
\subsection{Qiskit-Aer (Noisy Simulation Backend)}
\subsection{Qiskit-IBM-Runtime (Real Hardware)}
\subsection{FastAPI}
\subsection{WebSockets for Real-Time Streaming}
\subsection{SQLite for Run Persistence}
\subsection{React with Vite}
\subsection{Tailwind CSS}

% =================== 5. APPLICATION ARCHITECTURE =============================
\newpage
\section{Application Architecture}
\label{sec:architecture}

\subsection{The C4 Model}

\subsubsection{Level 1: System Context}
\subsubsection{Level 2: Containers}
\subsubsection{Level 3: Components}
\subsubsection{Level 4: Code}

\subsection{Database Schema}

% =================== 6. IMPLEMENTATION =======================================
\newpage
\section{Implementation and Features}
\label{sec:implementation}

\subsection{User Flow Diagram}
\subsection{Configuration Panel}
\subsection{The BB84 Protocol Engine}
\subsection{Eve Eavesdropping Simulation}
\subsection{IBM Quantum Hardware Integration}
\subsection{Real-Time WebSocket Streaming}
\subsection{PhotonGrid Visualization}
\subsection{Key Distillation Pipeline}
\subsection{QBER Historical Chart}
\subsection{Educational Panel}
\subsection{Simulation History and Persistence}

% =================== 7. EXPERIMENTAL RESULTS =================================
\newpage
\section{Experimental Results}
\label{sec:results}

\subsection{Methodology}
\subsection{Simulator vs.\ Real IBM Hardware}
\subsection{QBER Distribution under Different Eve Modes}
\subsection{Impact of Noise Parameters on Key Yield}

% =================== 8. CONCLUSIONS ==========================================
\newpage
\section{Conclusions}
\label{sec:conclusions}

\subsection{Future Development Possibilities}
\subsection{General Achievements and Conclusions}

% =================== 9. LIST OF FIGURES ======================================
\newpage
\listoffigures

% =================== 10. BIBLIOGRAPHY ========================================
\newpage
\begin{thebibliography}{99}

  \bibitem{bb84}
    C.~H. Bennett and G.~Brassard,
    \textit{Quantum Cryptography: Public Key Distribution and Coin Tossing},
    Proceedings of the IEEE International Conference on Computers, Systems and
    Signal Processing, Bangalore, India, pp.~175--179, 1984.

  \bibitem{shor1994}
    P.~W. Shor,
    \textit{Algorithms for Quantum Computation: Discrete Logarithms and
    Factoring},
    Proceedings of the 35th Annual Symposium on Foundations of Computer
    Science (FOCS), pp.~124--134, IEEE, 1994.

  \bibitem{mosca2018}
    M.~Mosca,
    \textit{Cybersecurity in an Era with Quantum Computers: Will We Be Ready?},
    IEEE Security \& Privacy, 16(5), pp.~38--41, 2018.

  \bibitem{wootters1982}
    W.~K. Wootters and W.~H. Zurek,
    \textit{A Single Quantum Cannot Be Cloned},
    Nature, 299, pp.~802--803, 1982.

  \bibitem{mayers2001}
    D.~Mayers,
    \textit{Unconditional Security in Quantum Cryptography},
    Journal of the ACM, 48(3), pp.~351--406, 2001.

  \bibitem{loChau1999}
    H.-K. Lo and H.~F. Chau,
    \textit{Unconditional Security of Quantum Key Distribution over
    Arbitrarily Long Distances},
    Science, 283(5410), pp.~2050--2056, 1999.

  \bibitem{qiskit2024}
    Qiskit contributors,
    \textit{Qiskit: An Open-source Framework for Quantum Computing},
    \url{https://qiskit.org}, 2024.

  \bibitem{grover1996}
    L.~K. Grover,
    \textit{A Fast Quantum Mechanical Algorithm for Database Search},
    Proceedings of the 28th Annual ACM Symposium on Theory of Computing
    (STOC), pp.~212--219, 1996.

  \bibitem{shorPreskill2000}
    P.~W. Shor and J.~Preskill,
    \textit{Simple Proof of Security of the BB84 Quantum Key Distribution
    Protocol},
    Physical Review Letters, 85(2), pp.~441--444, 2000.

  \bibitem{bennett1995pa}
    C.~H. Bennett, G.~Brassard, C.~Cr\'{e}peau, and U.~M. Maurer,
    \textit{Generalized Privacy Amplification},
    IEEE Transactions on Information Theory, 41(6), pp.~1915--1923, 1995.

  \bibitem{brassard1994}
    G.~Brassard and L.~Salvail,
    \textit{Secret-Key Reconciliation by Public Discussion},
    Advances in Cryptology --- EUROCRYPT 1993, LNCS 765,
    pp.~410--423, Springer, 1994.

  \bibitem{ITU-T-G652}
    International Telecommunication Union,
    \textit{Characteristics of a Single-Mode Optical Fibre and Cable
    (ITU-T Recommendation G.652)}, 2016.

  \bibitem{gllp2004}
    D.~Gottesman, H.-K. Lo, N.~L\"utkenhaus, and J.~Preskill,
    \textit{Security of Quantum Key Distribution with Imperfect Devices},
    Quantum Information and Computation, 4(5), pp.~325--360, 2004.

  \bibitem{lucamarini2018}
    M.~Lucamarini, Z.~L. Yuan, J.~F. Dynes, and A.~J. Shields,
    \textit{Overcoming the Rate--Distance Limit of Quantum Key Distribution
    without Quantum Repeaters},
    Nature, 557(7705), pp.~400--403, 2018.

  \bibitem{lo2005decoy}
    H.-K. Lo, X.~Ma, and K.~Chen,
    \textit{Decoy State Quantum Key Distribution},
    Physical Review Letters, 94, 230504, 2005.

  \bibitem{wang2005decoy}
    X.-B. Wang,
    \textit{Beating the Photon-Number-Splitting Attack in Practical Quantum
    Cryptography},
    Physical Review Letters, 94, 230503, 2005.

  \bibitem{martinez2015cascade}
    J.~Mart\'{i}nez-Mateo, C.~Pacher, M.~Peev, A.~Ciurana, and V.~Mart\'{i}n,
    \textit{Demystifying the Information Reconciliation Protocol Cascade},
    Quantum Information and Computation, 15(5\&6), pp.~453--477, 2015.

  \bibitem{breiman2001}
    L.~Breiman,
    \textit{Random Forests},
    Machine Learning, 45(1), pp.~5--32, 2001.

  \bibitem{xu2024dlqkd}
    Y.~Xu et al.,
    \textit{Deep-Learning-Based Continuous Attacks on Continuous-Variable
    Quantum Key Distribution Protocols},
    arXiv:2408.12571, 2024.

  \bibitem{zhao2025rlqkd}
    S.~Zhao et al.,
    \textit{Deep Reinforcement Learning and Variational Autoencoders for
    Enhanced Quantum Key Distribution Security},
    Results in Engineering, 2025.

  \bibitem{b92}
    C.~H. Bennett,
    \textit{Quantum Cryptography Using Any Two Nonorthogonal States},
    Physical Review Letters, 68(21), pp.~3121--3124, 1992.

  \bibitem{sarg04}
    V.~Scarani, A.~Ac\'{i}n, G.~Ribordy, and N.~Gisin,
    \textit{Quantum Cryptography Protocols Robust against Photon Number
    Splitting Attacks for Weak Laser Pulse Implementations},
    Physical Review Letters, 92(5), 057901, 2004.

  \bibitem{ekert1991}
    A.~K. Ekert,
    \textit{Quantum Cryptography Based on Bell's Theorem},
    Physical Review Letters, 67(6), pp.~661--663, 1991.

  \bibitem{hochreiter1997}
    S.~Hochreiter and J.~Schmidhuber,
    \textit{Long Short-Term Memory},
    Neural Computation, 9(8), pp.~1735--1780, 1997.

  \bibitem{tomamichel2012}
    M.~Tomamichel, C.~C.~W. Lim, N.~Gisin, and R.~Renner,
    \textit{Tight Finite-Key Analysis for Quantum Cryptography},
    Nature Communications, 3, 634, 2012.

\end{thebibliography}

\end{document}
